\scanpage{17}
\S--- \uline{Commentaires sur la définition d'une cohomologie $p$-adique en car. $p>0$}

$X_0$ schéma sur corps parfait $k$ de car. $p>0$, $S=\unsure{\operatorname{Spec}|\operatorname{Spf}}W(k)$. Alors
\[
H^*(X_0/S,\mathrm{cris};\underline O_{X_0\,\mathrm{cris}})
\]
est un $W$-module gradué. \deleted{\unsure{L'espoir}} Mon espoir initial était qu'il soit \deleted{un} ``le'' bon candidat pour une cohomologie $p$-adique. Un premier \unsure{support} \deleted{\unsure{(moral)}} en faveur, c'est que lorsque $X_0$ \inserted{est lisse sur $k$, et pour simplifier,} se remonte en $X$ \deleted{propre et}\inserted{projectif et} lisse sur $\unsure{k|W}$, alors d'après l'exposé précédent, on trouve aisément
\[
H^*((X_0\deleted{\unsure{/S}})_{\mathrm{cris}},\underline O_{\deleted{\unsure{X_0\,\mathrm{cris}}}})
=\varprojlim_n H^*((X_n/\unsure{k_n|W_n})_{\mathrm{cris}},O)
\]
et j'espérais qu'à la limite, on trouverait (peut-être modulo torsion) la $\overline{\mathrm m}$ chose que $\varprojlim H^*_{\mathrm{DR}}(X_n/\unsure{k_n|W_n})$.
\inserted{$\simeq H^*_{\mathrm{DR}}(X/S)$.}
(Il y \deleted{\unsure{aurait}} a une flèche en tous cas !).
C'est faux malheureusement. Par exemple si $X_0/k$ est un schéma abélien ``ordinaire'' i.e. qui du point de vue formel est un tore, on trouve que
\[
H^i((X_0/S)_{\mathrm{cris}},\underline O)=0\quad\text{si }i>\dim X_0\ !
\]
(Pourtant la $\varprojlim$ des coh. de De Rham est excellente !)
\deleted{\unsure{Le calcul de}} Quand on regarde dans le cas
\margin{16}
